% META: {"id": "TNSI-BDD-COURS-01", "chapitre": "TNSI-BASES-DE-DONNEES", "type_objet": "cours", "capacites_codes": ["C1"], "sources_inspiration": [], "mode_creation": "ex_nihilo", "status": "approved", "capacites": ["TNSI-BASES-DE-DONNEES-C1"]}
\section{Le modèle relationnel}

% ============================================================
% STRATE 1 — Concepts du modele relationnel
% ============================================================

\definition[1]{
Le \textbf{modèle relationnel} organise les données en \textbf{relations} (tables). Une relation est constituée de \textbf{tuples} (lignes), chacun décrivant une occurrence, et d'\textbf{attributs} (colonnes), chacun défini sur un \textbf{domaine} (l'ensemble des valeurs possibles, par exemple entier ou texte). Le \textbf{schéma} d'une relation liste le nom de la relation et le nom de chacun de ses attributs.
}

\margeAppui{
On note le schéma d'une relation \texttt{Livre(id, titre, auteur, annee)} : \texttt{Livre} est le nom de la relation, \texttt{id}, \texttt{titre}, \texttt{auteur}, \texttt{annee} sont ses attributs.
}

\definition[2]{
Une \textbf{clé primaire} est un attribut (ou groupe d'attributs) dont la valeur identifie de façon \textbf{unique} chaque tuple d'une relation. Une \textbf{clé étrangère} est un attribut d'une relation dont la valeur fait référence à la clé primaire d'une autre relation : c'est ainsi que deux relations sont mises en lien.
}

\textbf{Exemple.} Une bibliothèque gère trois relations liées entre elles :
\begin{sql}
Livre(id, titre, auteur, annee)
Usager(id, nom, prenom)
Emprunt(id, id_livre, id_usager, date_emprunt, date_retour)
\end{sql}
Dans \texttt{Emprunt}, \texttt{id\_livre} est une clé étrangère vers \texttt{Livre.id}, et \texttt{id\_usager} une clé étrangère vers \texttt{Usager.id} : chaque emprunt référence un livre et un usager précis, sans dupliquer leurs informations.

% BEGIN-VERIFY
% import sqlite3
% conn = sqlite3.connect(":memory:")
% cur = conn.cursor()
% cur.execute("CREATE TABLE Livre(id INTEGER PRIMARY KEY, titre TEXT, auteur TEXT, annee INTEGER)")
% cur.execute("CREATE TABLE Usager(id INTEGER PRIMARY KEY, nom TEXT, prenom TEXT)")
% cur.execute("""CREATE TABLE Emprunt(
%     id INTEGER PRIMARY KEY,
%     id_livre INTEGER REFERENCES Livre(id),
%     id_usager INTEGER REFERENCES Usager(id),
%     date_emprunt TEXT,
%     date_retour TEXT
% )""")
% cur.execute("INSERT INTO Livre VALUES (1, 'Le Petit Prince', 'Saint-Exupery', 1943)")
% cur.execute("INSERT INTO Usager VALUES (1, 'Diallo', 'Awa')")
% cur.execute("INSERT INTO Emprunt VALUES (1, 1, 1, '2026-01-10', NULL)")
% conn.commit()
% cur.execute("SELECT titre FROM Livre JOIN Emprunt ON Livre.id = Emprunt.id_livre WHERE Emprunt.id_usager = 1")
% assert cur.fetchone() == ("Le Petit Prince",)
% END-VERIFY

% ============================================================
% STRATE 2 — Structure vs contenu, anomalies
% ============================================================

\propriete[Structure et contenu]{
La \textbf{structure} d'une base de données (son schéma : relations, attributs, clés, contraintes) est stable dans le temps. Le \textbf{contenu} (les tuples effectivement stockés) évolue à chaque insertion, modification ou suppression. Modifier le contenu ne change jamais la structure ; modifier la structure (ajouter une colonne, par exemple) laisse le contenu existant intact autant que possible.
}

\erreurFrequente{
\textbf{Confondre relation et fichier plat.} Stocker toutes les informations d'une bibliothèque dans une seule table \texttt{Emprunt(id, titre\_livre, auteur, nom\_usager, prenom\_usager, ...)} \textbf{duplique} le titre et l'auteur à chaque emprunt du même livre : c'est une \textbf{anomalie de schéma}. Une modification (correction d'une faute dans un titre) doit alors être répercutée sur toutes les lignes concernées, au risque d'incohérences. Séparer \texttt{Livre}, \texttt{Usager} et \texttt{Emprunt} en trois relations liées par des clés étrangères élimine cette duplication.
}

\approfondissement{
Ce principe de séparation pour éviter les anomalies porte un nom en théorie des bases de données : la \textbf{normalisation} (formes normales 1NF, 2NF, 3NF). Le programme de Terminale ne l'exige pas formellement, mais en comprendre l'intuition (une information ne doit être stockée qu'à un seul endroit) aide à concevoir des schémas robustes.
}
